La validación funcional del sistema se apoya en tres pruebas
\textit{end-to-end} que recorren los flujos críticos sin saltarse capas.
La prueba \texttt{test\_lifecycle\_e2e} parte de la ingesta de páginas
escaneadas y verifica todas las transiciones de la
\texttt{ExamInstance} hasta su archivo, comprobando que las incidencias
se registran y resuelven según las reglas del
\cref{SEC:DES-LIFECYCLE}. La prueba \texttt{test\_grading\_flow} valida
la asignación de correctores, la introducción de anotaciones, el
cálculo de la calificación por rúbrica y la transición a
\texttt{PUBLISHED}. La prueba \texttt{test\_export} verifica la
exportación de calificaciones a \acs{csv} con el formato configurable
descrito en \cref{SS:DES-OPS-ORG}. En conjunto, las tres pruebas
verifican el cumplimiento simultáneo de los requisitos
\cref{rf:instances,rf:annotations,rf:grading,rf:review} sobre datos
generados de forma realista.

Como complemento a las pruebas automatizadas, el sistema dispone de
tres comandos de demostración invocables desde el \texttt{Makefile} que
generan artefactos visualmente inspeccionables: \texttt{make demo-qr}
ejecuta el \textit{pipeline} completo de instrumentación de un \acs{pdf}
con códigos \acs{qr} y posterior decodificación;
\texttt{make demo-watermark} produce páginas con la marca de agua
aplicada en sus cinco estilos (\texttt{grid}, \texttt{diagonal},
\texttt{center\_stamp}, \texttt{bottom\_right}, \texttt{cross}) sobre
\acs{pdf} e imagen \textit{raster}; \texttt{make demo-ingestion} genera
un \textit{dataset} sintético de instancias y páginas para
\textit{benchmarks} de ingesta a escala. Estos comandos sirven como
verificación manual reproducible de las decisiones de diseño descritas
en \cref{SEC:DES-INGEST,SEC:DES-SECURITY}.
